% Part: computability
% Chapter: machines-computations
% Section: configuration

\documentclass[../../../include/open-logic-section]{subfiles}

\begin{document}

\olfileid[tr]{cmp}{tur}{con}
\olsection{Konfigürasyonlar ve Hesaplar}

\begin{explain}
Daha önce çift sayı makinesinin konfigürasyonlarını adım adım izlediğimizi
anımsayın. Şerit, okuma-yazma kafası ve Turing makinesi programından oluşan
hayalî düzenek, bir Turing makinesi hesabını görselleştirmenin sezgisel bir
yoludur. Biçimsel olarak bir Turing makinesinin belirli bir girdideki hesabını
\emph{konfigürasyonlar} dizisi olarak tanımlayabiliriz. Konfigürasyonun
kendisi, hesapta belirli bir andaki şerit içeriğine karşılık gelen simgeler
dizisi, okuma-yazma kafasının konumunu belirten bir sayı ve bir durumdur.
Bunları kullanarak $M$ Turing makinesinin belirli bir girdide ne hesapladığını
tanımlayabiliriz.
\end{explain}

\begin{defn}[Konfigürasyon]
$M=\tuple{Q,\Sigma,q_0,\delta}$ Turing makinesinin bir
\emph{konfigürasyonu}, aşağıdaki koşulları sağlayan
$\tuple{C,m,q}$ üçlüsüdür:
\begin{enumerate}
\item $C\in\Sigma^*$, $\Sigma$ simgelerinden oluşan sonlu bir dizidir.
\item $m\in\Nat$, $m<\len{C}$ olan bir sayıdır.
\item $q\in Q$.
\end{enumerate}
Sezgisel olarak $C$ dizisi şeridin içeriğidir; en soldaki kareden son boş
olmayan ya da daha önce ziyaret edilmiş kareye kadarki bütün karelerin
simgelerini verir. $m$, okuma-yazma kafasının taradığı karenin numarasıdır;
en soldaki kare $0$ ile başlar. $q$ ise makinenin o andaki durumudur.
\end{defn}

\begin{explain}
Bir Turing makinesinin olası girdisi, genellikle bir sayıyı belirli biçimde
kodlayan bir simgeler dizisidir. Turing makinesinin başlangıç konfigürasyonu,
makineyi bu girdi üzerinde çalıştırmaya başladığımız konfigürasyondur: Şerit,
hemen ardından sağdaki karelere yazılmış girdinin geldiği şerit sonu işaretini
içerir; okuma-yazma kafası girdinin en soldaki karesini, yani sol uç
işaretinin sağındaki kareyi tarar; düzenek belirtilmiş başlangıç durumu
$q_0$'dadır.
\end{explain}

\begin{defn}[Başlangıç konfigürasyonu]
$M$'nin $I\in\Sigma^*$ girdisi için \emph{başlangıç konfigürasyonu}
\[
  \tuple{\TMendtape\frown I,1,q_0}
\]
üçlüsüdür.
\end{defn}

\begin{explain}
$\frown$ simgesi \emph{art arda birleştirme}yi gösterir; girdi dizgesi sol
uç işaretinin hemen sağında başlar.
\end{explain}

\begin{defn}
Bir $\tuple{C,m,q}$ konfigürasyonunun $M$'ye göre
\emph{tek adımda $\tuple{C',m',q'}$ konfigürasyonunu verdiğini} ancak ve
ancak aşağıdaki koşullarda söyleriz:
\begin{enumerate}
\item $C$'nin $m$. simgesi $\sigma$'dır.
\item $M$'nin yönerge kümesi
  $\delta(q,\sigma)=\tuple{q',\sigma',D}$ değerini belirtir.
\item $C'$ dizisinin $m$. simgesi $\sigma'$'dür.
\item
  \begin{enumerate}
  \item $D=L$ ve $m>0$ ise $m'=m-1$; aksi durumda $m'=0$, ya da
  \item $D=R$ ve $m'=m+1$, ya da
  \item $D=N$ ve $m'=m$.
  \end{enumerate}
\item $m'=\len{C}$ ise $\len{C'}=\len{C}+1$ ve $C'$ dizisinin $m'$.
  simgesi $\TMblank$'dir. Aksi durumda $\len{C'}=\len{C}$.
\item $i<\len{C}$ ve $i\neq m$ olan her $i$ için $C'(i)=C(i)$.
\end{enumerate}
\end{defn}

\begin{defn}\ollabel{defn:run-output}
$M$'nin $I$ girdisindeki bir \emph{çalışması}, $M$ konfigürasyonlarından
oluşan öyle bir $C_i$ dizisidir ki $C_0$, $M$'nin $I$ için başlangıç
konfigürasyonu ve her $C_i$, tek adımda $C_{i+1}$ konfigürasyonunu verir.

$C_k=\tuple{C,m,q}$, $C$'nin $m$. simgesi $\sigma$ ve
$\delta(q,\sigma)$ tanımsızsa $M$'nin \emph{$I$ girdisinde $k$ adım sonra
durduğunu} söyleriz. Bu durumda $M$'nin $I$ girdisindeki \emph{çıktısı}
$O$'dur; burada $O$, sonunda $\TMblank$ bulunmayan ve bazı $j\in\Nat$ için
$C=\TMendtape\concat O\concat\TMblank^j$ olacak bir simgeler dizgesidir.
($\TMblank^j$, $j$ tane $\TMblank$ simgesinden oluşan dizidir.)
\end{defn}

\begin{explain}
Bu tanıma göre $M$'nin çıktısı $O$ her zaman $\TMblank$ dışında bir simgeyle
biter. $k$ anında en soldaki $\TMendtape$ dışında bütün şerit $\TMblank$ ile
doluysa $O$ boş dizgedir.
\end{explain}

\end{document}
